\documentclass[11pt,a4paper]{article}
\usepackage[T1]{fontenc}
\usepackage[utf8]{inputenc}
\usepackage[english]{babel}
\usepackage{amsmath,amssymb}
\usepackage[margin=2cm]{geometry}
\setlength{\parindent}{0pt}
\setlength{\emergencystretch}{3em}
\allowdisplaybreaks

\title{About \texttt{false\_book\_systems.tex}}
\author{}
\date{}
\begin{document}
\maketitle
This note explains what the test data in \texttt{false\_book\_systems.tex} is used for. That file is input for the engine in \texttt{main.c}; it is read by the test \texttt{test\_linear\_systems.c} in this folder.

\begin{enumerate}
\item \textbf{Purpose.} The file holds 34 false statements about the set of solutions of a system of linear equations. They are anti-regression data: the engine must declare each of them false, with a certificate that the independent checker accepts, and must never prove one.

\item \textbf{How they were made.} Each is a statement of \texttt{book\_systems.tex} (the sources are listed in \texttt{book\_systems\_about.tex}) with its result changed by little, in one of the ways in which a solution written by hand goes wrong:
\begin{itemize}
\item one component of a solution changed, or its sign (statements 1, 2, 9, 12, 13, 18, 24, 29, 31);
\item a vector multiplied by a parameter with a wrong component or a wrong sign (6, 10, 15, 17, 21, 34);
\item a parameter forgotten: one solution, or a smaller family, in place of all of them (5, 11, 16, 23, 28);
\item a parameter too many: a set larger than the set of solutions (25);
\item a system with solutions said to have none (4, 8, 26, 33), and a system without solutions said to have some, or given one (3, 7, 14, 22, 27);
\item a solution said not to be one (30), and the right set said to be different (32);
\item a result as it is printed: statement 19 is the general solution printed for exercise 92 of Vietri, whose last component has the wrong sign, and statement 20 the particular solution as it is misprinted there.
\end{itemize}

\item \textbf{A third computation.} Every statement was recomputed by the independent oracle \texttt{tools/fuzz/fuzz\_linear\_algebra.py systems-data}, with ranks computed by minors: it found all 34 false.

\item \textbf{What is being tested.} That a wrong result is refuted and not merely left unproved, in each of the ways of chapter 33 of \texttt{main.c}: the written vector does not satisfy an equation; a written vector does not solve the homogeneous system; the system has a solution that the written set does not contain, with the row that separates it; the system has a solution, or has none. And that the certificate of each refutation is accepted by the checker, rejected with the claim reversed, and no longer accepted with a witness removed or altered.

\item \textbf{Format expected by the test.} The same as \texttt{book\_systems.tex}.

\item \textbf{How to run it.} From the top folder of the repository:
\begin{verbatim}
cc -std=c11 -O2 -Wall -Wextra -Werror -o /tmp/test_linear_systems \
    tests/linear_algebra/test_linear_systems.c && /tmp/test_linear_systems
\end{verbatim}

\item \textbf{Expected result.} The line \texttt{PASS: 34 of 34 statements \dots{} are declared false with a certificate that the checker accepts, rejects with the claim reversed, and no longer accepts with a witness removed or altered}.
\end{enumerate}
\end{document}
